\section{核心代码}

\subsection{问题一：功率平衡计算}

\begin{verbatim}
import numpy as np

def solve_problem1():
    # 设备参数 (36吨/日产能)
    P_EHA_rated = 10 + 10 + 0.75  # 20.75 MW
    
    # 读取标幺功率曲线并转换为实际功率
    P_load = load_pu * 6.0  # MW
    P_wind = wind_pu * 40.0  # MW
    P_pv = pv_pu * 64.0  # MW
    
    # 功率平衡
    P_net = P_wind + P_pv - P_EHA_rated - P_load
    P_buy = np.maximum(-P_net, 0)
    P_sell = np.maximum(P_net, 0)
    
    # 电量计算
    E_total = np.sum(P_EHA_rated + P_load)
    E_RE = np.sum(P_wind + P_pv)
    E_buy, E_sell = np.sum(P_buy), np.sum(P_sell)
    
    # 绿电指标
    R1 = (E_total - E_sell - E_buy) / E_RE
    R2 = (E_RE - E_sell) / E_total
    R3 = E_sell / E_RE
    
    # 吨氨成本
    C_RE = np.sum(P_wind*0.15 + P_pv*0.12) * 1000
    C_OM = (10*0.1 + 10*0.15 + 0.75*0.002) * 24 * 1000
    C_buy = np.sum(P_buy * prices) * 1000
    C_sell_rev = E_sell * 0.3779 * 1000
    C_ton = (C_RE + C_OM + C_buy - C_sell_rev) / 36.0
    return C_ton, R1, R2, R3
\end{verbatim}

\subsection{问题二：贪心排序优化}

\begin{verbatim}
def optimize_discrete(N_on, P_wind, P_pv, P_load, prices):
    P_EHA = 41.5  # MW (72吨/日产能)
    P_RE = P_wind + P_pv
    marginal_cost = np.zeros(24)
    
    for t in range(24):
        # 停机时净成本
        surplus_off = P_RE[t] - P_load[t]
        cost_off = -surplus_off*0.3779*1000 if surplus_off>=0 \
                   else (-surplus_off)*prices[t]*1000
        # 开机时净成本
        surplus_on = P_RE[t] - P_load[t] - P_EHA
        cost_on = -surplus_on*0.3779*1000 if surplus_on>=0 \
                  else (-surplus_on)*prices[t]*1000
        # 运维成本
        OM = (20*0.1 + 20*0.15 + 1.5*0.002) * 1000
        marginal_cost[t] = cost_on + OM - cost_off
    
    # 选择边际成本最低的N_on个时段
    schedule = np.zeros(24)
    schedule[np.argsort(marginal_cost)[:N_on]] = 1
    return schedule
\end{verbatim}

\subsection{问题三：线性规划求解}

\begin{verbatim}
from scipy.optimize import linprog

def optimize_continuous(M_target, P_wind, P_pv, P_load, prices):
    P_EHA = 41.5
    P_net = P_wind + P_pv - P_load
    n = 24
    # 决策变量: [alpha(24), b(24), s(24)]
    c = np.zeros(3*n)
    c_om = (20*0.1+20*0.15+1.5*0.002)/P_EHA
    for t in range(n):
        c[t] = c_om * P_EHA * 1000 / M_target
        c[n+t] = prices[t] * 1000 / M_target
        c[2*n+t] = -0.3779 * 1000 / M_target
    
    # 等式约束: 产量 + 功率平衡
    A_eq = np.zeros((n+1, 3*n))
    b_eq = np.zeros(n+1)
    A_eq[0, :n] = 1.0
    b_eq[0] = M_target / 3.0
    for t in range(n):
        A_eq[1+t, t] = P_EHA
        A_eq[1+t, n+t] = -1.0
        A_eq[1+t, 2*n+t] = 1.0
        b_eq[1+t] = P_net[t]
    
    bounds = [(0.1,1.0)]*n + [(0,None)]*n + [(0,None)]*n
    res = linprog(c, A_eq=A_eq, b_eq=b_eq, bounds=bounds,
                  method='highs')
    alpha = res.x[:n]
    return alpha
\end{verbatim}

\subsection{问题四：离网储能优化}

\begin{verbatim}
def optimize_offgrid_storage(C_sto, P_wind, P_pv, P_load):
    P_EHA = 41.5
    eta_c, eta_d, sigma = 0.90, 0.90, 0.002
    P_net = P_wind + P_pv - P_load
    n = 24
    # 决策变量: [alpha(24), Pc(24), Pd(24), curtail(24)]
    nvar = 4 * n
    # 目标: 最大化产量 = 最小化 -sum(alpha)
    c = np.zeros(nvar)
    c[:n] = -3.0  # 最大化产量
    
    # 约束: 功率平衡 + 储能状态方程 + SOC边界
    # ... (LP构建省略，详见code/problem4.py)
    
    bounds = [(0.1,1.0)]*n + [(0,2*C_sto)]*n + \
             [(0,2*C_sto)]*n + [(0,None)]*n
    res = linprog(c, A_eq=A_eq, b_eq=b_eq,
                  A_ub=A_ub, b_ub=b_ub,
                  bounds=bounds, method='highs')
    return res.x[:n]  # 最优alpha
\end{verbatim}
